\chapter*{\textbf{Preface}}
\phantomsection
\addcontentsline{toc}{chapter}{Preface}

Around 1920, a journalist asked the British astrophysicist Arthur Eddington whether it was true that only three people in the world understood general relativity. Eddington paused for a moment and replied: ``I am trying to think who the third person is.'' The anecdote may or may not be historically accurate, but Eddington's claim 
certainly was not. Within months of Einstein publishing his field equations in 1915, exact solutions were already being found: Schwarzschild's black hole, the de Sitter and anti-de Sitter cosmologies, Friedmann's expanding universe. The theory proved far more accessible than its reputation suggested \cite{LlibreBertJanssen}.

That reputation, however, is not entirely undeserved. General 
relativity rests on tensor analysis and differential geometry, branches of mathematics that most physics students encounter only when they need them for this very purpose. Yet beneath the formalism lies a surprisingly simple physical idea: spacetime is a dynamical entity, shaped by matter and energy, and what we call gravity is nothing more than the curvature of that entity. As John Wheeler famously put it, ``matter tells spacetime how to curve, and spacetime tells matter how to move'' \cite{Misner1973}.

One hundred and ten years after the publication of the field 
equations, general relativity has passed every experimental test devised. Among its most striking predictions are gravitational waves, ripples in the fabric of spacetime produced by the most violent events in the universe. First detected in September 2015 by the LIGO Scientific Collaboration \cite{PrimeradeteccioGW150914}, they opened 
an entirely new observational window onto the cosmos, one based not on light but on the oscillation of spacetime itself. In the decade since, nearly two hundred compact binary mergers have been detected, and the field is entering what can only be described as its golden age.

This thesis sits at the intersection of two of the richest areas in modern gravitational-wave astronomy: strong gravitational lensing and sky localization. When a gravitational wave passes near a massive galaxy on its way to Earth, its path is deflected and it can arrive at the detector as multiple copies of the same signal, separated in time and amplitude. Identifying these lensed image pairs and analyzing them jointly opens the possibility of dramatically improving our knowledge of where on the sky the source lies. That is the central question this work sets out to address.